\scanpage{36}\mspage{35}
% @opening
\deleted{\unsure{Il faut que|figure} \unsure{l'on a|On a}} La condition est la suivante ; \deleted{\unsure{qu'on requiert}}
% @fini

\uline{Cont} \deleted{\unsure{4|2|3}}$3'$ (Facultative dans le cas non finitiste) \deleted{$\forall$} Pour tout $X\in M$, \deleted{la figure} \inserted{l'ens} $M_{\leq X}$ est fini. (i.e. les figures élémentaires sont finies)
% @dims

Par contre, on exige en tous cas que la \uuline{dimension combinatoire} de tout \deleted{\uline{Cont} 3 La} $X\in M$ soit finie (i.e. $\dim\text{ comb }M_{\leq X}$ finie), et plus \deleted{géné} précisément encore
% @mark
% slanted
\margin{\uuline{\unsure{devant|demain}}\hspace{2pt}\tikz[baseline]{\draw (0,-.2)--(0,.4) (.045,-.2)--(.045,.4) (.09,-.2)--(.09,.4);}}
% @sup
\[
\uline{\text{Cont}}\,3 \quad \text{Si } X\in M,\qquad
\underbrace{\mathop{\mathrm{Sup}}\limits_{X'\ll X}\,\mathrm{dim\,comb}\,M_{\leq X}}_{\text{\uline{dimension} (géom) de }X}<\infty
\]
% @place
On appelle \uuline{lieu} un \unsure{élémt|élément} de $M$ minimal pour $\ll$. On dénote par $\mathcal L$
% @subset
\[
\mathcal L\subset M
\]
% @endplace
l'ens des lieux.
% @shadow

Pour tout $X$, son \uuline{ombre} \inserted{(fermée)} sur $\mathcal L$ est l'ens $\mathcal L_{X}$ des lieux $x$ \deleted{$\unsure{\ll X}$} ``de $X$'' i.e. tels que $x\ll X$. On pose
% @sidelf
% slanted
\margin{$\mathcal L_F=\bigcup_{X\in F}\mathcal L_X$ s'appelle \uline{l'ombre de} \uline{la} figure \uline{$F$}}
% @lf
\[
\mathcal L_X^{\circ}=\mathcal L_X\setminus\bigcup_{Y<X}\mathcal L_Y
\]
% @axiom

\uline{Cont} \deleted{5}4 \quad a) $\forall X\in M$, on a $\mathcal L_X^{\circ}\neq\varnothing$

% @axiomb
b) Si $X,Y$ sont compatibles et distincts on a
% @disjoint
\[
\mathcal L_X^{\circ}\cap \mathcal L_Y^{\circ}=\varnothing
\]
